% !TeX root = ../main.tex
\section{Dự kiến doanh thu và chi phí trong giai đoạn đầu (12 tháng đầu)}
Các số liệu dự phóng tài chính trong năm đầu tiên được xây dựng dựa trên nguyên tắc thận trọng (Conservative approach).

\textbf{Cơ sở logic của các giả định tài chính:}
\begin{itemize}
    \item \textbf{Giá trị đơn thuê trung bình (AOV - Average Order Value):} Dựa trên khảo sát giá thuê thiết bị gaming phổ biến (30.000đ - 200.000đ/ngày), lấy trung bình là \textbf{50.000 VNĐ / 1 ngày} và hành vi thuê thường kéo dài 2-3 ngày dịp cuối tuần, AOV được ước tính ở mức an toàn là \textbf{150.000 VNĐ/đơn}.
    \item \textbf{Lãi gộp cơ sở (Gross Margin):} Mức tiền cọc trung bình được giả định khoảng \textbf{2.000.000 VNĐ}, tương ứng với giá trị của các thiết bị gaming phổ biến (bàn phím cơ, chuột gaming, tai nghe, tay cầm...) trên nền tảng. 
    \begin{itemize}
        \item Phí mở ví là \textbf{30.000 VNĐ / 1 tháng} và người thuê được sử dụng trong mức tiền cọc cho phép. Tỉ lệ phần trăm chia giữa Dịch vụ hỗ trợ tiền cọc (Deposit Advance Service) và Mutux là \textbf{50\%} vậy nên lợi nhuận thu được là \textbf{15.000 VNĐ / 1 tháng / 1 user}. Tỉ lệ người sử dụng ví là \textbf{70\%} người sử dụng, còn lại sẽ là người dùng tiền họ tự nạp.
        \item Doanh thu từ đơn thuê sẽ chiếm \textbf{30\%} trên giá trị đơn thuê. Giả sử với \textbf{AOV} = 150.000 VNĐ thì lợi nhuận mỗi đơn Mutux nhận được là \textbf{45.000 VNĐ}.
    \end{itemize}
\end{itemize}

\textbf{Tổng giá trị giao dịch trên nền tảng theo từng quý - GMV (Gross Merchandise Value):}

\begin{table}[htbp]
    \centering
    \caption{Tổng giá trị giao dịch theo từng quý (GMV)}
    \begin{tabular}{@{}lccccc@{}}
        \toprule
        \textbf{Giai đoạn} & \textbf{User ví/Tổng User} & \textbf{Số đơn} & \textbf{GD theo đơn} & \textbf{GD mở ví} & \textbf{Tổng GD Quý (GMV)} \\
        \midrule
        \textbf{Quý 1} (Thử nghiệm) & 140/200 & 300 & 45.000.000 & 4.200.000 & \textbf{49.200.000} \\
        \textbf{Quý 2} (Tăng trưởng) & 300/450 & 550 & 82.500.000 & 9.000.000 & \textbf{91.500.000} \\
        \textbf{Quý 3} (Chiếm lĩnh) & 700/1000 & 2.000 & 300.000.000 & 21.000.000 & \textbf{321.000.000} \\
        \textbf{Quý 4} (Mở rộng \& B2B) & 1400/2000 & 3.500 & 525.000.000 & 42.000.000 & \textbf{567.000.000} \\
        \midrule
        \textbf{TỔNG NĂM 1} & 1400/2000 & 6.350 & 952.500.000 & 76.200.000 & \textbf{1.028.700.000} \\
        \bottomrule
    \end{tabular}
\end{table}

\textbf{Doanh thu theo từng quý và cả năm sau khi chia tiền với các bên}

\begin{table}[htbp]
    \centering
    \caption{Doanh thu dự phóng năm 1}
    \begin{tabular}{@{}lrrrr@{}}
        \toprule
        \textbf{Giai đoạn} & \textbf{GMV} & \textbf{Doanh thu thuê} & \textbf{Doanh thu mở ví} & \textbf{Doanh thu tổng} \\
        \midrule
        \textbf{Quý 1} & 49.200.000 & 13.500.000 & 2.100.000 & \textbf{15.600.000} \\
        \textbf{Quý 2} & 91.500.000 & 24.750.000 & 4.500.000 & \textbf{29.250.000} \\
        \textbf{Quý 3} & 321.500.000 & 100.000.000 & 10.500.000 & \textbf{110.500.000} \\
        \textbf{Quý 4} & 567.000.000 & 157.500.000 & 21.000.000 & \textbf{178.500.000} \\
        \midrule
        \textbf{CẢ NĂM} & 1.028.700.000 & 285.750.000 & 38.100.000 & \textbf{323.850.000} \\
        \bottomrule
    \end{tabular}
\end{table}

\textbf{Cash Flow - dòng chảy tiền qua từng Quý}

\begin{table}[htbp]
    \centering
    \caption{Dòng chảy tiền (Cash Flow) các quý năm đầu}
    \begin{tabular}{@{}lrrrrr@{}}
        \toprule
        \textbf{Quarter} & \textbf{Thu (Tr VNĐ)} & \textbf{Chi (Tr VNĐ)} & \textbf{Net Burn} & \textbf{Cash Flow} & \textbf{Cash Drain} \\
        \midrule
        Q1 & 15 & 120 & -105 & -105 & -105 \\
        Q2 & 29 & 120 & -91 & -91 & -196 \\
        Q3 & 110 & 120 & -10 & -10 & -206 \\
        Q4 & 178 & 120 & +58 & +58 & -148 \\
        \bottomrule
    \end{tabular}
\end{table}

Bắt đầu từ năm thứ 2, các doanh thu tăng theo số lượng người dùng, số lượng đơn hàng. Các nguồn thu mới như nhận quảng cáo trên nền tảng, đẩy các bài đăng gaming gear, phí hoa hồng mua luôn sản phẩm sau khi thuê. Dù vậy doanh thu chính vẫn đến từ số lượng đơn thuê.

\textbf{Tổng kết tài chính Năm 1:} 
Tổng chi phí vận hành trong năm đầu tiên của Mutux được ước tính khoảng 480 triệu đồng. Tuy nhiên, doanh nghiệp dự kiến tạo ra khoảng 324 triệu đồng doanh thu trong cùng giai đoạn. Do đó, khoản thiếu hụt dòng tiền tích lũy (Funding Gap) vào khoảng 156 triệu đồng.

Để đảm bảo doanh nghiệp có đủ dòng tiền duy trì hoạt động liên tục, xử lý các rủi ro phát sinh và không bị gián đoạn trong giai đoạn đầu, Mutux dự kiến chuẩn bị nguồn vốn khoảng \textbf{250--300 triệu đồng}. Nguồn vốn này có thể đến từ vốn tự có của nhóm sáng lập, các chương trình hỗ trợ khởi nghiệp hoặc vòng gọi vốn Pre-seed. Chi tiết các thông số:
\begin{itemize}
    \item Burn rate: Ta có \textit{Gross Burn} hay chi phí mỗi tháng là 480/12 = 40 triệu, \textit{Net Burn} sau khi đã bù doanh thu theo từng quý đã được tính theo bảng trên.
    \item Vậy nên nhu cầu tìm kiếm nhà đầu tư để tăng hoặc có vốn làm runway trong ít nhất 2 năm tương đương \textbf{$\sim$300 triệu}. Với 2 năm đầu làm nền tảng thì Mutux có thể chạm đến điểm hòa vốn và bắt đầu lời. Từ đó bứt phá và chiếm lĩnh thị trường cho thuê gaming gear và phát triển thêm.
\end{itemize}

\textbf{Điểm hòa vốn (Break-even):} 
Đặt mục tiêu đạt được vào khoảng \textbf{đầu năm thứ 2 (tháng 13 - 15)}. Lộ trình này hoàn toàn phù hợp với vòng đời của một Startup Marketplace. Khi lượng người dùng thường xuyên đi vào ổn định ở năm thứ 2, chi phí thu hút khách hàng (CAC) sẽ giảm dần. Khi đó, mức doanh thu kỳ vọng sẽ đủ để bù đắp chi phí vận hành hàng tháng theo kịch bản tài chính của dự án (ước tính dao động khoảng 80 - 100 triệu VNĐ/tháng).
